\documentclass[10pt,a4paper]{amsart}
\usepackage[margin=19mm]{geometry}
\usepackage{amsmath,amssymb,mathtools}
\usepackage[scr=rsfs,cal=euler]{mathalfa}
\usepackage{xfrac}
\usepackage[unicode,hidelinks,bookmarks=false]{hyperref}
\newcommand{\ie}{i.e.\ }
\newcommand{\cf}{cf.\ }
\newcommand{\resp}{respectively\ }
\newcommand{\Subsection}{\subsection}
\providecommand{\nobreakdash}{\nobreak\hbox{-}}
\setlength{\emergencystretch}{2em}
\multlinegap0pt
\usepackage{amsthm}
\usepackage{bm}
\DeclareMathAlphabet\mathbfcal{OMS}{cmsy}{b}{n}
\renewcommand{\Re}{\operatorname{Re}}
\renewcommand{\Im}{\operatorname{Im}}
\newcommand{\dist}{\operatorname{dist}}
\newcommand{\LI}{{\mathbb T}}
\newcommand{\LJ}{{\overline {\mathbb T}}}
\newcommand{\LII}{{\LI\LI}}
\newcommand{\LIJ}{{\LI\LJ}}
\newcommand{\LJI}{{\LJ\LI}}
\newcommand{\LJJ}{{\LJ\LJ}}
\newtheorem{theorem}{Theorem}
\newtheorem{lemma}{Lemma}
\newtheorem{corollary}{Corollary}
\newtheorem{definition}{Definition}
\newtheorem{proposition}{Proposition}
\title[On truncations of hierarchical equations of motion]{On truncations of hierarchical equations of motion for finite-dimensional systems: spectral convergence of the truncated Liouvillians (Theorem 2)}
\author[Vasilii Vadimov]{Vasilii Vadimov}
\date{Source: arXiv:2604.22568v2 (CC BY 4.0)}
\begin{document}
\maketitle
\vspace{1em}
\noindent\textit{Source and licence.} On truncations of hierarchical equations of motion for finite-dimensional systems. Version arXiv:2604.22568v2 (CC BY 4.0), distributed under the Creative Commons Attribution 4.0 International licence, http://creativecommons.org/licenses/by/4.0/, as recorded in the arXiv OAI licence metadata for this exact version and shown on the abstract page https://arxiv.org/abs/2604.22568v2. This item is an editorial re-presentation of the paper's spectral convergence theorem together with the paper's own proof of it, the retained source-local chain it uses, and the single bibliography entry the retained text cites, transcribed from the paper's own bibliography. A full rights notice is delivered at licenses/arxiv-2604.22568-CC-BY-4.0.txt.
\noindent\textit{Editorial scope.} Counted unit: the paper's spectral convergence theorem for truncated Liouvillians, printed as Theorem 2 (Spectral convergence) (source0.tex lines 688--702), delivered together with the paper's complete original proof of it (source0.tex lines 704--856) as the single primary proof body. No mathematics is written, repaired or re-derived: statement and proof are the paper's own text line for line, in the paper's own notation ($\mathbfcal L$, $\mathbfcal L_{\LI_k}$, $\sigma_K$, $K_\pm$, $\mathbfcal R(z;\mathbfcal L)$, $\mathbfcal S_{\LI_k}$, $\LI_k$). Required context retained uncounted, because the counted proof is the paper's own assembly over it: the decay-ratio lemma (source0.tex lines 609--612), the discreteness corollary with its complete proof (lines 615--629), and the Schur-complement convergence lemma with its statement (lines 655--685). Every cross-reference inside the retained spans resolves inside the delivered document, and the only citation commands of the delivered text are the two of the retained corollary proof and counted proof, whose single entry is transcribed from the paper's own bibliography with the paper's own printed number (33). Printed numbers are the paper's own: the wrapper restores the paper's per-kind theorem counters and its global equation counter before the retained material, so the retained decay-ratio lemma prints as Lemma 3, the retained discreteness corollary as Corollary 2, the retained Schur-complement lemma as Lemma 4, the counted statement as Theorem 2, and the numbered displays of the retained spans carry the paper's own equation numbers (29) to (43). Omitted from this package and not redistributed: the paper's preamble and front matter as such (source0.tex lines 1--60 and the title block of lines 61--250, of which only the title and the author name are reproduced in the wrapper); the introduction and model section; the Gershgorin-type resolvent bound and block decomposition of the paper's Section 3; the stability section and the discussion, conclusions and appendices that follow; and the remaining 36 entries of the paper's own bibliography. The delivered document therefore carries a bibliography of exactly the one work its retained text cites, with the paper's own printed number, and no reference or citation command of the delivered text was rewritten or adapted. Printed slips kept literal, among them the paper's own wording and its own choice of symbols. Layout and class: the paper's own class is the publisher's SciPost class, which is present in the e-print and is neither shipped with this package nor used to build it; the delivered wrapper is the lane's pinned amsart editorial wrapper at 10pt on A4, to which this package adds the paper's own theorem-like declarations with their separate per-kind counters, its own calligraphic alphabet declaration and its own shorthand macros, all copied verbatim from the paper's source. No publisher class file, no publisher style file and no upstream script is redistributed. Assets: no retained span of this package contains a figure environment, a drawing picture or a graphics-inclusion command, no image or artwork file exists in or is redistributed with this package, and the manifest and the delivery evidence both carry an empty asset-dependency list. A full rights notice is delivered at licenses/arxiv-2604.22568-CC-BY-4.0.txt.
\setcounter{lemma}{2}
\begin{lemma}
    For any exhausting sequence of truncations \(\{\LI_k\}_{k=1}^\infty\), $\Gamma'_{\LI_k} \to +\infty$ and $\Gamma_{\LI_k} / \Gamma'_{\LI_k} \to 1$ as $k \to \infty$.
    \label{lemma:decay-ratio}
\end{lemma}

\setcounter{corollary}{1}
\begin{corollary}
    A bounded set \(\Omega \subset \mathbb C\) contains at most finitely many spectral values of \(\mathbfcal L\).
    \label{lemma:spectrum-discreteness}
\end{corollary}
\begin{proof}
    Choose an exhausting sequence of truncations. According to Lemma~\ref{lemma:decay-ratio}, $\Gamma_\LI \to +\infty$ along this sequence, so we can pick a truncation \(\LI\) such that
       $ -\Gamma_\LI < \inf_{z \in \Omega} \Re z$, 
    hence
    $
        \overline{\Omega}
        \subset
        \Omega\left(\mathbfcal L_\LJJ\right).
    $
    Therefore, \(\mathbfcal S_\LI(z)\) is holomorphic in a neighborhood of \(\overline{\Omega}\). Since \(\LI\) is finite, \(\det \mathbfcal S_\LI(z)\) is a scalar holomorphic function. Its zeros in the compact set \(\overline{\Omega}\) are isolated and hence finite in number~\cite{Ablowitz2003}. Since $\mathbfcal S_\LI(z)$ is the Schur complement of $z - \mathbfcal L$, these zeros are precisely the spectral values of \(\mathbfcal L\).
\end{proof}

\setcounter{lemma}{3}
\setcounter{equation}{28}
\begin{lemma}
    For a truncation $\LI$ such that $\Gamma_\LI > 0$ and a complex $z$ such that $\Re z>-\Gamma_\LI$, let
    \begin{equation}
        \varepsilon_\LI(z)
        :=
        \frac{
            C\sqrt{(\Upsilon_\LI+\gamma)(\Upsilon_\LI+\tilde\gamma)}
        }{\Gamma_\LI}
        \left[
            \frac{|z|}{\Re z+\Gamma_\LI}
            +
            \frac{\sqrt{C(\Gamma_\LI'+\gamma)}}{\Gamma_\LI'}
        \right].
        \label{eq:schur-error-majorant}
    \end{equation}
    Then
    \begin{equation}
        \left \| \mathbfcal S_\LI(z) - z + {\mathbfcal L}_\LI \right \|_1
        \leqslant
        \varepsilon_\LI(z).
        \label{eq:schur-error-bound}
    \end{equation}
    Moreover, for any bounded $\Omega \subset \mathbb C$ and any exhausting sequence of truncations $\{\LI_k\}_{k=1}^{\infty}$,
    \begin{equation}
        \lim_{k\to\infty}
        \sup_{z\in\Omega}
        \varepsilon_{\LI_k}(z)
        =0.
    \end{equation}
    \label{lemma:schur-convergence}
\end{lemma}

\setcounter{theorem}{1}
\begin{theorem}[Spectral convergence]
Let \(K \subset \mathbb C\) be a compact set. For each \(\lambda \in \sigma_K(\mathbfcal L) := \sigma(\mathbfcal L) \cap K\), choose a bounded open neighborhood \(V_\lambda\) with rectifiable boundary \(\partial V_\lambda\) such that
    $\overline V_\lambda \cap \sigma(\mathbfcal L) = \{\lambda\}$,
and let
    $K_- := K \setminus \bigcup_{\lambda \in \sigma_K(\mathbfcal L)} V_\lambda$.
    Then for any exhausting sequence of truncations $\{\LI_k\}_{k=1}^\infty$ there is $k_0$ such that for all $k > k_0$ the following statements hold:
\begin{enumerate}
    \item
        Each neighborhood \(V_\lambda\) contains at least one eigenvalue of \(\mathbfcal L_{\LI_k}\). Moreover, the number of eigenvalues of \(\mathbfcal L_{\LI_k}\) in \(V_\lambda\), counted with algebraic multiplicity, coincides with the algebraic multiplicity of \(\lambda\) as an eigenvalue of \(\mathbfcal L\).

    \item
        The operator \(\mathbfcal L_{\LI_k}\) has no eigenvalues in \(K_-\).
\end{enumerate}
\label{th:spectrum-convergence}
\end{theorem}

\begin{proof}
    Let
        $K_+
        :=
        K \cup \bigcup_{\lambda \in \sigma_K(\mathbfcal L)} \overline{V_\lambda}$.
    Since \(K\) is compact and \(\sigma_K(\mathbfcal L)\) is finite by Lemma~\ref{lemma:spectrum-discreteness}, the set \(K_+\) is compact.

    For each \(\lambda \in \sigma_K(\mathbfcal L)\), define
    \begin{equation}
        M_\lambda
        :=
        \max_{z \in \partial V_\lambda}
        \left \|
        \mathbfcal R(z;\mathbfcal L)
        \right \|_1 .
    \end{equation}
    This maximum exists because \(\partial V_\lambda\) is compact and disjoint from \(\sigma(\mathbfcal L)\). Likewise, define
    \begin{equation}
        M_-
        :=
        \max_{z \in K_-}
        \left \|
        \mathbfcal R(z;\mathbfcal L)
        \right \|_1 ,
    \end{equation}
    which is finite because \(K_-\) is a compact subset of \(\varrho(\mathbfcal L)\). Finally, set
    \begin{equation}
        M := \max\left(M_-, \max_{\lambda \in \sigma_K(\mathbfcal L)} M_\lambda\right).
    \end{equation}

    By Lemma~\ref{lemma:schur-convergence} there exists \(k_0\) such that for all $k > k_0$
    \begin{equation}
        \sup_{z \in K_+}
        \left \|
        \mathbfcal S_{\LI_k}(z) - z + \mathbfcal L_{\LI_k}
        \right \|_1
        <
        \frac{1}{M}.
        \label{eq:small-schur-error}
    \end{equation}

    We first prove the statement for a fixed \(\lambda \in \sigma(\mathbfcal L)\cap K\). For \(t \in [0,1]\), define
    \begin{equation}
        \mathbfcal F_{\LI_k}(z,t)
        :=
        t\left(z-\mathbfcal L_{\LI_k}\right)
        +
        (1-t)\mathbfcal S_{\LI_k}(z).
    \end{equation}
    Since
    \begin{equation}
        \mathbfcal F_{\LI_k}(z,t)
        =
        \left[
            \mathbfcal I
            -
            t\left(
                \mathbfcal S_{\LI_k}(z)-z+\mathbfcal L_{\LI_k}
            \right)
            \mathbfcal S_{\LI_k}(z)^{-1}
        \right]
    \mathbfcal S_{\LI_k}(z),
    \end{equation}
    it is enough to show that the bracket is invertible for all $t\in[0, 1]$ and all $z \in \partial V_\lambda$. By the Schur-complement representation,
        $\mathbfcal S_{\LI_k}(z)^{-1}
        =
        \mathbfcal R_{\LI_k \LI_k}(z;\mathbfcal L)
        $,
    and therefore
    \begin{equation}
        \left \|
        \mathbfcal S_{\LI_k}(z)^{-1}
        \right \|_1
        \leqslant
        \left \|
        \mathbfcal R(z;\mathbfcal L)
        \right \|_1
        \leqslant
        M
        \qquad
        \text{for } z \in \partial V_\lambda.
    \end{equation}
    Together with~\eqref{eq:small-schur-error}, this yields
    \begin{equation}
        \left \|
        t\left(
            \mathbfcal S_{\LI_k}(z)-z+\mathbfcal L_{\LI_k}
        \right)
        \mathbfcal S_{\LI_k}(z)^{-1}
        \right \|_1
        < 1,
    \end{equation}
    so \(\mathbfcal F_{\LI_k}(z,t)\) is invertible for all \(z \in \partial V_\lambda\) and all \(t \in [0,1]\).
    Since \(\LI_k\) is finite, \(\mathbfcal F_{\LI_k}(z,t)\) is matrix-valued.      Let us denote
    $f_t(z) := \det \mathbfcal F_{\LI_k}(z,t)$.
    For each fixed \(t\), the function \(f_t(z)\) is holomorphic in \(V_\lambda\) and continuous on \(\partial V_\lambda\). Since $\mathbfcal F_{\LI_k}(z, t)$ is invertible,
    $f_t(z) \neq 0$
    for all  $z \in \partial V_\lambda$, $t \in [0,1]$.
    Hence the logarithmic derivative \(f_t'(z)/f_t(z)\) is holomorphic in a neighborhood of \(\partial V_\lambda\), and the argument principle~\cite{Ablowitz2003} gives
    \begin{equation}
        N_\lambda(t)
        :=
        \frac{1}{2\pi i}
        \oint_{\partial V_\lambda}
        \frac{f_t'(z)}{f_t(z)}\,\mathrm dz,
        \label{eq:Nt-argument-principle}
    \end{equation}
    where \(N_\lambda(t)\) is exactly the number of zeros of \(f_t\) in \(V_\lambda\), counted with multiplicity. For \(t=0\), $N_\lambda(t)$ is equal to the number of the spectral values of \(\mathbfcal L\) in \(V_\lambda\). For \(t=1\), $N_\lambda(t)$ gives the number of the eigenvalues of \(\mathbfcal L_\LI\) in \(V_\lambda\).

    We show that \(N_\lambda(t)\) is continuous in \(t\). Since \(\partial V_\lambda \times [0,1]\) is compact and \(f_t(z)\) is continuous in $(z,t)$, the function \(|f_t(z)|\) attains a strictly positive minimum on \(\partial V_\lambda \times [0,1]\), therefore $1 / f_t(z)$ is also continuous in $(z, t)$. Moreover, because \(\mathbfcal F_{\LI_k}(z,t)\) is continuous in \((z,t)\) on \(\partial V_\lambda \times [0,1]\) and holomorphic in \(z\), the derivative \(f_t'(z)\) depends continuously on \((z,t)\) on \(\partial V_\lambda \times [0,1]\). Hence, the logarithmic derivative ${f_t'(z)}/{f_t(z)}$ is continuous on \(\partial V_\lambda \times [0,1]\). It follows that the integral in~\eqref{eq:Nt-argument-principle} depends continuously on \(t\), so \(N_\lambda(t)\) is a continuous function on \([0,1]\). On the other hand, \(N_\lambda(t)\) is integer-valued by the argument principle. Being both continuous and integer-valued on the connected set \([0,1]\), it must be constant, therefore $N_\lambda(0) = N_\lambda(1)$.

    We now prove the second statement. For \(z \in K_-\), we can express
    \begin{equation}
        z-\mathbfcal L_{\LI_k}
        =
        \left[
            \mathbfcal I
            -
            \left(
                \mathbfcal S_{\LI_k}(z)-z+\mathbfcal L_{\LI_k}
            \right)
            \mathbfcal S_{\LI_k}(z)^{-1}
        \right]
    \mathbfcal S_{\LI_k}(z).
    \end{equation}
    Since \(K_- \subset \varrho(\mathbfcal L)\), we have
    \begin{equation}
        \left \|
        \mathbfcal S_{\LI_k}(z)^{-1}
        \right \|_1
        \leqslant
        \left \|
        \mathbfcal R(z;\mathbfcal L)
        \right \|_1
        \leqslant
        M
        \qquad
        \text{for all } z \in K_-.
    \end{equation}
    Hence, by~\eqref{eq:small-schur-error},
    \begin{equation}
        \left \|
        \left(
            \mathbfcal S_{\LI_k}(z)-z+\mathbfcal L_{\LI_k}
        \right)
        \mathbfcal S_{\LI_k}(z)^{-1}
        \right \|_1
        < 1
        \qquad
        \text{for all } z \in K_-.
    \end{equation}
    Therefore, \(z-\mathbfcal L_{\LI_k}\) is invertible for all \(z \in K_-\), and \(\mathbfcal L_{\LI_k}\) has no eigenvalues in \(K_-\).
\end{proof}

\begin{thebibliography}{99}
\bibitem[33]{Ablowitz2003} M.~J. Ablowitz and A.~S. Fokas, \newblock \emph{Complex variables: introduction and applications}, \newblock Cambridge university press (2003).
\end{thebibliography}
\end{document}
